\documentclass[../../../include/open-logic-section]{subfiles}

\begin{document}

\olfileid{sth}{story}{predicative}	

\olsection{直谓与非直谓}

罗素集合 $R$ 是通过 $\Setabs{x}{x \notin x}$ 来定义的。更完整地说，$R$ 将是包含且仅包含所有那些非自属集合的集合。因此，在定义 $R$ 时，我们对本应包含 $R$ 的论域（如果 $R$ 存在的话）进行了量化。

这是一个\emph{非直谓}的定义。更一般地说，我们可以说，一个定义是非直谓的，当且仅当它对一个包含被定义对象自身的论域进行量化。
	
在悖论出现之后，Whitehead、Russell、Poincaré 和 Weyl 将此类非直谓定义斥为“恶性循环的”：

\begin{quote}
	对所要避免的悖论进行分析表明，它们都源于某种恶性循环。这里所说的恶性循环之所以产生，是由于假定一个对象的汇集可能包含这样的成员：这些成员只有借助该汇集作为整体才能定义[\ldots \textparagraph]

	使我们能够避免不合法全体的原则可以陈述如下：“凡涉及一个汇集的\emph{所有}成员者，其自身不得是该汇集的一员”；或者反过来说：“如果某个汇集一旦构成了一个整体，它就将包含只能借助该整体来定义的成员，那么该汇集就根本没有整体。”我们将称此为“恶性循环原则”，因为它使我们能够避免因假定不合法全体而涉及的恶性循环。
	\citep[p.~37]{WhiteheadRussell1910}
\end{quote}

如果我们追随他们，依据恶性循环原则进行拒斥，那么我们就可能尝试用类似下面的方案来取代（参见 \olref[sth][story][rus]{sec}）那灾难性的朴素概括模式：

\begin{defish}
\emph{直谓概括。} 对于每一个仅对集合进行量化的公式 $\phi$：集合$^\prime$ $\Setabs{x}{\phi(x)}$ 存在。
\end{defish}

只要集合$^{\prime}$不是集合，就不会产生矛盾。

不幸的是，直谓概括并不是那么\emph{包罗}。毕竟，它向我们引入了新的实体，即集合$^\prime$。因此我们将不得不考虑那些对集合$^\prime$进行量化的公式。如果它们总是生成一个集合$^\prime$，那么罗素悖论将再次出现，只需要考虑所有非自属的集合$^\prime$所构成的集合$^\prime$即可。所以，沿着同样的思路，我们必须说，一个对集合$^\prime$进行量化的公式生成的是相应的集合$^{\prime\prime}$。而后我们将需要集合$^{\prime\prime\prime}$、集合$^{\prime\prime\prime\prime}$，等等。为了避免撇号泛滥，更方便的做法是将它们视为集合$_0$、集合$_1$、集合$_2$、集合$_3$、集合$_4$……这就为我们通向（简单）类型论提供了一条路径。

对于这样一种理论，存在着一些显而易见的反对意见（尽管这些意见是否\emph{压倒性的}并不显然）。简言之：由此产生的理论使用起来很繁琐；它在设定不同种类的对象时过于铺张；并且最终来看，并不清楚非直谓定义到底有多\emph{糟}。
	
为了阐明最后一点，请看 \citeauthor{Ramsey1925} 的这一评论：

\begin{quote}
	我们可以称一个人在某个群体中是最高的，从而借助一个他自身也是其成员的全体来识别他，而这并不包含任何恶性循环。 \citep{Ramsey1925}
\end{quote}

Ramsey 的观点是，“群体中最高的人”\emph{确实}是一个非直谓定义；但它显然是完全正当的。

有人可能会回应说，在这个案例中，我们可以通过\emph{直谓}的方式找出那个最高的人。例如，也许我们可以直接指出那个人。那么，对非直谓定义的反对显然就需要限定于那些“只能”非直谓地指称的实体。但即便如此，我们还需要更多的解释，来说明为什么这种“本质上的非直谓性”会如此糟糕。\footnote{更多讨论，参见 \citet{Linnebo2010}。}

诚然，如果我们希望自己的定义能为我们提供类似“创造”对象的配方，那么非直谓定义绝对是坏消息。因为，给定一个非直谓定义，人们将真正陷入恶性循环：要创造那个被非直谓地指定的对象，首先需要创造所有对象（包括那个被非直谓地指定的对象本身），因为该对象是依据所有对象来指定的；所以，在创造出该对象自身之前，就需要先把它创造出来。但这同样只构成针对本质上非直谓地指定的集合的严重反对意见，前提是我们认为集合是我们\emph{创造}的东西。而我们（很可能）并不这样认为。

因此——无论好坏——后来变得常见的方法并不涉及对（非）直谓性采取强硬立场。相反，它采用的是如今被视为累积迭代的进路。最终，这将允许我们把集合分层为“阶段”——有点类似于直谓进路将实体分层为集合$_0$、集合$_1$、集合$_2$……——但我们不会假定它们之间存在种类的差异。

\end{document}